%% 
%% Copyright 2007, 2008, 2009 Elsevier Ltd
%% 
%% This file is part of the 'Elsarticle Bundle'.
%% ---------------------------------------------
%% 
%% It may be distributed under the conditions of the LaTeX Project Public
%% License, either version 1.2 of this license or (at your option) any
%% later version.  The latest version of this license is in
%%    http://www.latex-project.org/lppl.txt
%% and version 1.2 or later is part of all distributions of LaTeX
%% version 1999/12/01 or later.
%% 
%% The list of all files belonging to the 'Elsarticle Bundle' is
%% given in the file `manifest.txt'.
%% 

%% Template article for Elsevier's document class `elsarticle'
%% with numbered style bibliographic references
%% SP 2008/03/01

% Final journal-style: two columns



\UseRawInputEncoding
\documentclass[final,5p,times]{elsarticle}


\usepackage{microtype}

% Revisions / change markup
\usepackage{xcolor}
\newcommand{\revision}[1]{\textcolor{red}{#1}}


%% Use the option review to obtain double line spacing
%\documentclass[review,14pt,a4paper]{elsarticle}

%% For including figures, graphicx.sty has been loaded in
%% elsarticle.cls. If you prefer to use the old commands
%% please give \usepackage{epsfig}

%% The amssymb package provides various useful mathematical symbols
\usepackage{amssymb}
\usepackage{hyperref}
\setlength{\parindent}{0pt}
%% The amsthm package provides extended theorem environments
%% \usepackage{amsthm}

%% The lineno packages adds line numbers. Start line numbering with
%% \begin{linenumbers}, end it with \end{linenumbers}. Or switch it on
%% for the whole article with \linenumbers.
\usepackage{lineno}
%\linenumbers
\journal{SoftwareX}


\biboptions{sort&compress}

\usepackage{amsmath}


\begin{document}
	\renewcommand{\labelenumii}{\arabic{enumi}.\arabic{enumii}}
	
	\begin{frontmatter}
		
		%% Title, authors and addresses
		
		%% use the tnoteref command within \title for footnotes;
		%% use the tnotetext command for theassociated footnote;
		%% use the fnref command within \author or \address for footnotes;
		%% use the fntext command for theassociated footnote;
		%% use the corref command within \author for corresponding author footnotes;
		%% use the cortext command for theassociated footnote;
		%% use the ead command for the email address,
		%% and the form \ead[url] for the home page:
		%% \title{Title\tnoteref{label1}}
		%% \tnotetext[label1]{}
		%% \author{Name\corref{cor1}\fnref{label2}}
		%% \ead{email address}
		%% \ead[url]{home page}
		%% \fntext[label2]{}
		%% \cortext[cor1]{}
		%% \address{Address\fnref{label3}}
		%% \fntext[label3]{}
		
		\title{DelaDect: Optical delamination detection in fiber-reinfoced polymers}
		
		%% use optional labels to link authors explicitly to addresses:
		%% \author[label1,label2]{}
		%% \address[label1]{}
		%% \address[label2]{}
		
		\author[label1]{Vasco D. C. Pires}
		\author[label2]{Matthias Rettl}
		\author[label1]{Martin Pletz \corref{cor1}}
		\author{Clara Schuecker  \fnref{label1}}
		\address[label1]{Chair of Designing Plastics and Composite Materials, Department of Polymer Engineering and Science, Technical University of Leoben, Austria}
		\address[label2]{Materials Center Leoben Forschung GmbH, Austria}
		
		\cortext[cor1]{Corresponding author. E-mail address: vasco.castro-pires@unileoben.ac.at}
		
		\begin{abstract}
			%% Text of abstract 
			Deladect is a tool for automated delamination detection in composite laminates obtained from image sequences.
			The package features a modular structure, enabling straightforward adaptation to different imaging setups and materials. By providing a reproducible workflow, DelaDect offers a standardized approach to delamination detection and quantification for both static and fatigue loading conditions. Distributed as an open-source Python package, it can further be extended to applications requiring detection and quantification of areas.
		\end{abstract}
		
		\begin{keyword}
			%% keywords here, in the form: keyword \sep keyword
			Delamination detection \sep Crack detection \sep Fiber-reinforced polymers \sep Python
			
			%% PACS codes here, in the form: \PACS code \sep code
			
			%% MSC codes here, in the form: \MSC code \sep code
			%% or \MSC[2008] code \sep code (2000 is the default)
			
		\end{keyword}
		
	\end{frontmatter}
	
	%\linenumbers
	
	\section*{Metadata}
	\phantomsection
	\label{sec:metadata}
	
	\begin{table*}[t]
		\centering
		\begin{tabular}{|l|p{6.5cm}|p{6.5cm}|}
			\hline
			\textbf{Nr.} & \textbf{Code metadata description} & \textbf{Metadata} \\
			\hline
			C1 & Current code version & v1.0.1 \\
			\hline
			C2 & Permanent link to the code used for this version &
			\url{https://github.com/VascoPires/DelaDect/commit/FULL_COMMIT_HASH} \\
			\hline
			C3 & Permanent link to reproducible capsule & Not available \\
			\hline
			C4 & Legal code license &
			GNU Affero General Public License v3.0 only (AGPL-3.0-only) \\
			\hline
			C5 & Code versioning system used & Git \\
			\hline
			C6 & Software languages, tools, and services used &
			Python; CrackDect; NumPy; SciPy; pandas; Matplotlib; scikit-image;
			Pillow; SQLAlchemy; Numba and psutil \\
			\hline
			C7 & Compilation requirements, operating environments, and dependencies &
			Python 3.10 \\
			\hline
			C8 & Developer documentation/manual &
			\url{https://deladect.readthedocs.io/en/latest/} \\
			\hline
			C9 & Support email for questions &
			vasco.castro-pires@unileoben.ac.at \\
			\hline
		\end{tabular}
		\caption{Code metadata.}
		\label{tab:code-metadata}
	\end{table*}
	
	\section{Motivation and significance}
	\label{loc:motivation_and_significance}
	Delamination in fiber-reinforced composite laminates is a major challenge in research and industry. It is one of the most critical failure modes in composites and can severely compromise the structural integrity of components \cite{wisnomRoleDelaminationFailure2012, demoraisEvaluationInitiationCriteria2006, huangReviewDelaminationDamage2023}. This phenomenon is often triggered by other damage mechanisms such as matrix cracks which usually appear in the most misaligned plies with the loading direction (also known as of \textit{off-axis cracks}) and are considered the first signs of damage with a measurable decrease in laminate stiffness \cite{talrejaStiffnessPropertiesComposite1986, nairnMicrocrackingMicrocrackinducedDelamination1992, bartley-choDamageAccumulationQuasiisotropic1998, tohgoFatigueBehaviorCFRP2006}. \revision{Although these cracks do not cause catastrophic failure, the local stress concentrations generated at the crack tips promote micro-debonding between plies which then evolve into higher delamination areas \cite{stinchcombNondestructiveEvaluationDamage1986, berthelotTransverseCrackingDelamination2003}. This type of delamination, which nucleates at the crack tips, is often called crack-induced delamination or \textit{diffuse delamination}. Furthermore, these effects are exacerbated close to the specimen edges, caused by additional out-of-plane stresses due to the existence of a free surface (edge effect) \cite{Hayashi1968AnalyticalSO, soniAnalysisFreeEdge1987, xuInteractionMatrixCracking1994, burhanPredictingFreeEdge2025}. This type of delamination which nucelates at the edges of the specimen is called \textit{edge delamination}.}
	
	The previously mentioned damage mechanisms can be categorized into stages based on their chronological appearance. Matrix cracks are categorized as Stage I damage, then the onset and growth of delamination are classified as Stage II and finally Stage III is the final failure of the component either caused by delamination, fiber-breakage or a mix of both. These damage mechanisms are schematically shown in \autoref{fig:damage_plot_3d_paper.png}. 
	
	\begin{figure*}
		\centering
		\includegraphics[width=\textwidth]{Files/damage_plot_3D_paper.png}
		\caption{Damage evolution in a composite laminate. Stage I corresponds to the onset of matrix cracking, which progresses until saturation is reached. In Stage II, these cracks act as nucleation sites for delamination driven by interlaminar stresses and, finally, Stage III is characterized by the continued growth of these damage mechanisms, ultimately leading to the fracture of the laminate.\label{fig:damage_plot_3d_paper.png}}
	\end{figure*}
	
	\revision{This progression from matrix cracking to delamination under static and cyclic loading has been studied experimentally in the past decades \cite{kimExperimentalAnalyticalStudies1984, nairnInitiationGrowthDelaminations1992, takedaInitiationGrowthDelamination1994, gludAutomatedCountingOffaxis2016, carraroCharacterisationAnalysisTransverse2019, liUnfoldingEarlyFatigue2020, quinlanAssessmentFatigueDamage2023}.  In these experiments, various imaging techniques are used to capture the onset and growth of damage, ranging from simple optical methods \cite{berthelotTransverseCrackingDelamination2003, garrettMultipleTransverseFracture1977, gludAutomatedCountingOffaxis2016, mehdikhaniMultiscaleDigitalImage2018, drvodericComparingCrackDensity2023, olesenMatrixCraCSAutomatedTracking2024} to advanced evaluation techniques like X-ray computed tomography (CT) \cite{hufenbachTestDeviceDamage2012, sketAutomaticQuantificationMatrix2014, garceaXrayComputedTomography2018, aiStressFieldDamage2021}. One of the most simple optical method is transmitted white light imaging (TWLI). In this setup, a light source is placed behind the specimen and a camera records the transmitted light on the opposite side, such that internal damage scatters the light and appears as dark areas in the images \cite{garrettMultipleTransverseFracture1977, berthelotTransverseCrackingDelamination2003, gludAutomatedCountingOffaxis2016, quinlanAssessmentFatigueDamage2023}. A clear limitation of this method is being only suitable for translucence materials such as glass fiber reinforced polymers (GFRPs). However, because of its simplicity, it is still commonly used to study cracking and delamination in composite laminates more broadly. Nonetheless, the interpretation and quantification of damage in these experiments is still performed manually, for crack counting \cite{addenStiffnessDegradationFatigue2010, bogerImprovementFatigueLife2010, quaresiminDamageEvolutionCyclic2014, kalteremidouInfluenceDifferentLayUp2021} and measuring delaminated areas/lengths \cite{takedaInitiationGrowthDelamination1994, kalteremidouInfluenceDifferentLayUp2021}. }
	
	\revision{
		To reduce this bottleneck, there has been a growing interest in automated image processing and detection of damage in composites. Glud et al. \cite{gludAutomatedCountingOffaxis2016} introduced an automated method for counting off-axis cracks from individual images, while later approaches used the complete image sequence to improve detection and tracking of cracks \cite{gludAutomatedCountingOffaxis2016, benderEffectVariableAmplitude2021}. These developments led to tools such as \emph{CrackDect} \cite{drvodericVersion02CrackDect2022} and \emph{MatrixCraCS}  \cite{olesenMatrixCraCSAutomatedTracking2024}. Moreover, some of these ideas, can be extended to images obtained from X-ray micro-computed tomography ($\mu$CT) \cite{sketAutomaticQuantificationMatrix2014, salgadoAutomatedImageProcessing2026}. Despite differences in the implementation, these methods generally rely on a combination of image normalization, directional filtering, intensity thresholding, and geometric criteria to detect cracks and track their evolution \cite{gludAutomatedCountingOffaxis2016, benderEffectVariableAmplitude2021, drvodericVersion02CrackDect2022, olesenMatrixCraCSAutomatedTracking2024, JustKochGude2022}. Despite the high degree of maturity for automated crack detection from images, automated delamination detection remains largely undeveloped.}
	
	\revision{
		For delamination detection, Carraro et al. \cite{carraroCharacterisationAnalysisTransverse2019}, proposed a methodology based on a front and back lightning systen. Under front illumination, the delamination regions appeared as white strips and were processed using an in-house automated Matlab\textsuperscript{\textregistered} routine based on image filtering and binarization. The resulting binary images were used to quantify delamination. Adam et al. \cite{adamSelectiveLaserIllumination2018} proposed a selective laser-illumination (SLI) method that allowed matrix cracks and delamination to be illuminated separately and used an automated tool capable of quantifying cracks and delaminations in both SLI and TWI obtained images. However, the automated tool was restricted to cross-ply laminates. Bak and Lindgaard \cite{bakMethodAutomatedDigital2020} developed an automated image-based method for tracking the position and shape of delamination fronts in GFRP specimens during delamination resistance fracture tests. In this case the delaminated area detected has a predictable and continuous shape. Together, these studies show that automated optical delamination evaluation is possible. However, fully automating the evaluation of delamination remains challenging since the detected regions can exhibit irregular morphologies \cite{shenExperimentalInvestigationDamage2017, adamSelectiveLaserIllumination2018}. Moreover, the available implementations are directed to specific experimental conditions and are neither openly available nor reproducible.}
	
	\revision{To address this gap, here \emph{DelaDect} is proposed, an open-source tool for detecting, segmenting, and quantifying delamination in composite laminates. Since delamination corresponds to Stage II of the damage evolution shown in \autoref{fig:damage_plot_3d_paper.png}, the framework also incorporates the analysis of the preceding Stage I matrix cracks for evaluation of diffuse delamination, through functionalities adapted from \emph{CrackDect} \cite{drvodericVersion02CrackDect2022}. The benefit of this tool is that it allows the fully automated categorization of Stage I and Stage II damage states in a composite using a fairly inexpensive method. Such outputs are particularly useful as either inputs or validation of numerical models. Although developed for TWLI images of GFRP laminates, its modular processing workflow can be adapted to other materials and imaging techniques, provided that damaged and undamaged regions have sufficient image contrast.}
	
	\section{Software description}
	\label{loc:software_description}
	
	\revision{\emph{DelaDect} is a Python package available on PyPI \cite{pypi2026} for image-based crack and delamination detection. Its crack, diffuse-delamination, and edge-delamination pipelines can be used independently or combined according to the user's requirements.}
	
	
	
	
	\subsection{Software architecture}
	\label{loc:software_description.software_architecture}
	DelaDect can be organized into the following subsystems (see \autoref{fig:flowchart_deladect.png}):
	\begin{itemize}
		\item \textbf{data structure and image handling}: shift correction (motion compensation), memory management and image loading;
		\item \textbf{crack detection}: \revision{generates frame-by-frame crack lists};
		\item \textbf{delamination detection}: normalization, filtering and delamination detection in each frame.
	\end{itemize}
	
	These subsystems are coordinated by the \emph{Specimen} class, which centralizes metadata and image stacks in a single data structure. \revision{The interactions between these subsystems are shown in \autoref{fig:flowchart_deladect.png} and described in the following sections.} \emph{DelaDect} is built on the commonly used Python libraries NumPy \cite{harrisArrayProgrammingNumPy2020}, SciPy \cite{virtanenSciPy10Fundamental2020}, Pandas \cite{mckinney2010pandas},  and scikit-image \cite{van2014scikit}, with visualization via Matplotlib \cite{hunterMatplotlib2DGraphics2007}. 
	
	\begin{figure*}
		\centering
		\includegraphics[width=\textwidth]{Files/flowchart_deladect.png}
		\caption{ Structure of the package's main subsystems: \textit{data structure and image handling}, \textit{crack detection} and \textit{delamination detection}. \revision{The first module handles image paths, preprocessing steps, and the operation of the \emph{ImageStack} class.} Crack detection is the module where cracks are detected from the image stack based on the CrackDect. In the delamination detection module, delamination areas around the previously detected cracks and specimen's edge are evaluated. An additional I/O module for saving/loading specimens and previous results is part of the package but not shown here.\label{fig:flowchart_deladect.png}}
	\end{figure*}
	\subsubsection{Data structure and image handling}
	\label{loc:software_description.software_architecture.data_structure_and_image_handling}
	The first step before any evaluation is the preprocessing of the images. During mechanical testing, sequence of images exhibit specimen drift within the field of view, and strain introduces local shear and rotation. As a result, the same feature may appear at different positions throughout the test, which can lead to the same damage being detected multiple times and being erroneously added as a new entry. \revision{So, some form of image correction is required.}  For this purpose, the software provides a script which includes a dedicated GUI for correcting translation, rotation and strain-induced distortions with multiple anchor points. \revision{For user convenience, these functionalities are included in the tool, although third-party software may also be used for image preparation.}
	
	After this motion compensation, the image sequences are loaded through the inherited \texttt{ImageStack} class family from \emph{CrackDect}. Two implementations are available: \texttt{ImageStack} keeps the complete sequence loaded in memory and offers fast access to each frame, while \texttt{ImageStackSQL} saves the image stack in a lightweight SQLite cache to minimize RAM consumption. At initiation, the \emph{DelaDect}  package estimates the memory footprint based on the image size and number of frames, and compares it with a user-configurable limit. Based on that, it automatically selects the most suitable backend. Users can override this selection if desired. Once the stack of images is defined, it is accessible for crack or delamination detection. For more operations of the \texttt{ImageStack} class, the reader is referred to the technical documentation of \cite{drvodericVersion02CrackDect2022}.
	\subsubsection{Crack analysis}
	\label{loc:software_description.software_architecture.crack_analysis}
	\revision{The crack detection layer uses \emph{CrackDect} \cite{drvodericVersion02CrackDect2022}, with the method of Bender et al.\ \cite{benderEffectVariableAmplitude2021} selected for its robustness and lower number of false detections, although it requires the complete image sequence. For individual frames, the simpler method of Glud et al. \cite{gludAutomatedCountingOffaxis2016} can be used, although it is more sensitive to noise.}
	
	\subsubsection{Delamination analysis}
	\label{loc:software_description.software_architecture.delamination_analysis}
	
	As mentioned before, the implemented delamination methodologies are not as developed as crack detection and also not as well documented. That said, the core ideas behind some of the well developed crack detection algorithms can be transferred to delamination detection, provided that the feature representation and processing steps are adapted to the delamination morphology. As shown in \autoref{fig:flowchart_deladect.png}, \emph{DelaDect} offers two \revision{methodologies to detect both types of delaminations.} Both follow the same overall logic, but  differ in the feature representation they target and therefore in the specific processing steps used. \revision{Furthermore, the diffuse-delamination pipeline requires prior crack detection, whereas the edge-delamination pipeline operates directly on the images. In both cases, the general delamination detection algorithm consists of the following operations (also shown in \autoref{fig:workflow_edge.png}):}
	
	\begin{enumerate}
		\item Minimum history: \revision{applies a global} minimum to the image stack. This enforces a monotonic darkening through time, since the images should only get darker due to damage, this makes the method more robust to bright illumination flashes or sensor noise through the test \cite{benderEffectVariableAmplitude2021};
		\item Normalization: \revision{the frame under analysis $I_i (x,y)$ is divided by a reference image $I_r (x.y)$. As long as both images are matched (thus the importance of shift correction), a change in intensity corresponds to a change in the relative damage state between both images \cite{benderEffectVariableAmplitude2021, JustKochGude2022}.} 
		\begin{equation}
			I^n_i (x,y) = \frac{I_i(x,y)}{I_r(x,y)},
			\label{eq:normalization}
		\end{equation}
		\revision{Here, the reference $I_r(x,y)$ can be either the first image in the stack, referred to as the \emph{static} approach, or the median of a set of recent frames, referred to as the \emph{rolling median} approach (see Step 2 of \autoref{fig:workflow_edge.png}). The rolling median clearly separates newly developed damage from damage already present in the specimen, however a higher number of frames is recomended.}
		\item Max/min filtering: \revision{morphological closing (maximum filter followed by a minimum filter for the same window). This suppresses thin, dark structures such as cracks and preserves broader delamination regions. The size of the window should be related with the average crack width.}
		\revision{
			\item Sharpening and Gaussian smoothing:
			\begin{itemize}
				\item \emph{Unsharp masking} [scikit-image] \cite{van2014scikit}: widens the intensity gap between delaminated and intact regions;;
				\item \emph{Gaussian smoothing} [SciPy] \cite{virtanenSciPy10Fundamental2020}: a Gaussian filter is applied to suppress noise and to smooth out intensity variations in the delamination areas.
		\end{itemize}}
		\item Constant scaling:  maps intensities to a fixed \revision{normalized} bounded range before thresholding, so threshold behaviour stays stable across frames \cite{bakMethodAutomatedDigital2020}.
		\item Thresholding: Converts the processed grayscale image into a binary candidate mask based on the k-means method \revision{($\text{k}=2$)} \cite{bakMethodAutomatedDigital2020, olesenMatrixCraCSAutomatedTracking2024} with Otsu threshold \cite{otsuThresholdSelectionMethod1979} \revision{as a fallback if convergence is not reached.} 
		\item Morphological closing: fills small holes and bridges narrow gaps left after thresholding in the candidate mask \cite{bakMethodAutomatedDigital2020}.
		\item Accumulation: updates each frame mask by running union with previous detections, enforcing physically consistent monotonic growth and reducing frame-to-frame flicker.
	\end{enumerate}
	
	\revision{After these processing steps, the delaminated area is computed from the number of pixels in the final binary mask.}
	
	\begin{figure}
		\centering
		\includegraphics[width=0.45\textwidth]{Files/workflow_edge.png}
		\caption{ Delamination detection steps common to both edge and diffuse delamination, shown for edge delamination on a $[\pm30/90]_s$ specimen.\label{fig:workflow_edge.png}}
	\end{figure}
	
	\medskip
	\noindent\textbf{Edge delamination}\par
	\smallskip
	\noindent
	Edge-delamination detection targets connected delaminated regions near the specimen edges. \revision{After applying the processing steps shown in \autoref{fig:workflow_edge.png}, newly detected regions are compared with the accumulated damage mask, and only regions connected to the specimen edges are retained as edge delamination. The tool also supports edge delamination analysis in multiple interfaces, although this functionality is beyond the scope of the present paper and is described in the software documentation.} 
	
	\medskip
	\noindent\textbf{Diffuse delamination}\par
	\smallskip
	\noindent
	Diffuse delamination targets crack-induced interior damage. It uses the crack catalogue to define local regions of interest (ROIs) around each detected crack, which are rotated and aligned with the local crack orientation. The common backbone operations are then applied within each ROI (Steps 1-7), and each ROI-level mask is mapped back into the full-frame coordinate system. The specimen-wide mask is assembled as shown in \autoref{fig:diffuse_workflow.png}.
	
	When edge and diffuse delamination are run in combined mode, overlaps are resolved with an edge-priority rule to prevent double counting: pixels claimed by both masks are assigned to edge delamination. Regardless of mode, the detected delamination area is constrained to be non-decreasing over time. For each frame, the area is obtained by pixel counting of the binary mask(s), yielding an area in $\text{pixels}^2$, which can then be converted to $\text{mm}^2$.
	
	\begin{figure*}[h!]
		\centering
		\includegraphics[width=0.8\textwidth]{Files/diffuse_workflow.png}
		\caption{ Diffuse delamination workflow used in \textit{DelaDect}. Cracks are detected first; segmentation is then performed in ROIs around each crack, followed by assembly into a specimen-wide mask.\label{fig:diffuse_workflow.png}}
	\end{figure*}
	
	\section{Illustrative examples}
	\label{loc:software_description.illustrative_examples}
	\revision{For one specimen, two code snipets are provided. The first code snippet shows how to set up the Specimen class and performs crack evaluation and the second code snippet is part of the delamination detection of both  diffuse and edge delamination.}
	
	\subsection{Crack evaluation}
	\label{loc:software_description.illustrative_examples.crack_evaluation}
	The example specimen is a cross-ply laminate with 0$^{\circ}$ and 90$^{\circ}$ plies, so crack detection is run separately for each ply orientation. \revision{In this example, a \emph{Specimen} object is first initialized with the image path and scale. The $0^{\circ}$ and $90^{\circ}$ plies and their interface are then defined, after which crack detection is performed for both ply orientations using the \texttt{crack\_analysis} function.}
	
	\includegraphics[width=\linewidth]{Files/crack_eval.pdf}
	The cracks detected by the example code snipet \revision{and the corresponding crack densities for both plies} are shown in \autoref{fig:results.png}.
	
	\subsection{Delamination evaluation (edge \& diffuse)}
	\label{loc:software_description.illustrative_examples.delamination_evaluation_(edge_+_diffuse)}
	For delamination detection, the same specimen is used. In addition to the ply definitions, an interface must be defined at the location where delamination should be evaluated. 
	
	\includegraphics[width=\linewidth]{Files/delam_eval.pdf}
	
	This produces the results shown in \autoref{fig:results.png}. \revision{The overlays in \autoref{fig:results.png}(a) are generated directly by the tool, while panels (b) and (c) were plotted using Matplotlib and PyVista, respectively. Other plotting tools can be used instead, as long as they can read DelaDect's standard output formats (see documentation).} 
	
	\begin{figure*}[h!]
		\centering
		\includegraphics[width=\textwidth]{Files/results.png}
		\caption{ \revision{Outputs generated by the code snippet: (a) detected cracks and delaminated regions overlaid on the corresponding raw images; (b) evolution of crack density and delaminated area; and (c) three-dimensional representation of the crack positions and delamination masks within the specimen, rendered using PyVista. For interpretation of the colours, the reader is referred to the online version of this article.}}
		\label{fig:results.png}
	\end{figure*}
	
	
	As with any optical damage detection method, result quality depends strongly on lighting conditions, focus, and contrast. Furthermore, poor crack resolution reduces the reliability of diffuse ROIs and can directly affect the diffuse delamination evaluation. So, it is recommended to use checkpoints for verifying the crack evaluation and preprocessing parameters before running an extensive delamination study. In practice, robust results typically require some initial parameter adjustment and user judgment. \revision{An in-depth explanation of each parameter and its role is provided in the software's documentation.  However, very broadly, parameters used in the previous example play the following roles:
		\begin{itemize}
			\item \texttt{window\_edge} and \texttt{window\_diffuse}: control how strongly crack-like features are suppressed (Step 3). Larger windows suppress more cracks but risk suppressing delamination along with it.
			\item \texttt{seed\_ratio}: the fraction of rows, starting from the specimen edge, used as seeds for delamination. Delamination can only be initiated starting from these rows (not shown in \autoref{fig:workflow_edge.png}, but available in the documentation).
			\item \texttt{diffuse\_dx} and \texttt{diffuse\_dy}: control the size of the rotated ROI within diffuse delamination is detected. A value of 1-2 times the average crack width is recommended as a starting point.
		\end{itemize}
		Other parameters/options are also available. For example, the methodology chosen to perform the normalization (Step 2), \texttt{gaussian\_filters}, which governs the Gaussian smoothing (Step 4), and \texttt{post\_threshold\_closing\_px}, which governs the post-threshold morphological closing (Step 7). Here they are left at their defaults to keep this illustrative example simple.}
	
	
	\section{Impact}
	\label{loc:impact}
	Quantifying delamination together with matrix cracking is vital for advancing the understanding of damage mechanisms in laminate composites and for improving damage models under both static and fatigue conditions. \textit{DelaDect} addresses this need by providing delamination quantification, allowing researchers to shift from qualitative observations to reproducible and quantitative analysis. The framework is capable of mapping delamination areas with minimal user input with a modular design which allows crack detection, diffuse delamination, and edge delamination to be detected independently or in combination. This structure also ensures compatibility with current workflows, for example, the crack detection can be substituted for another crack detection tool such as \cite{nicolAutomatedCrackDetection2023} or \cite{olesenMatrixCraCSAutomatedTracking2024} as long as the crack lists have the required format for \textit{DelaDect}. In addition, the filtering, and thresholding steps can be swapped or rewritten for other novel methods in the future without the need of major code rewrites. The package is scalable to thousands of high-resolution frames, making it suitable for long-duration fatigue experiments on standard hardware, and although tailored to composite delamination, its architecture can be extended to other area quantification tasks where areas must be tracked over time.
	\section{Conclusions}
	\label{loc:conclusions}
	In this work the modular delamination detection tool \textit{Deladect} is introduced. \textit{DelaDect} closes a practical gap in the damage detection of fiber reinforced composites, by providing a consistent bridge from crack  detection to delamination detection while still remaining modular and flexible. The resulting workflow is designed to be reproducible across test campaigns and scalable to large image sequences. The package is intended to help accelerate the development of composite damage models for both static and fatigue conditions. \revision{Its open-source nature also allows users to modify or extend the workflow with custom image-processing operations, supporting its adaptation and further development by the research community.}
	
	% \begin{enumerate}
		
		% \item \textbf{Minimum history.} A running minimum is maintained over all frames seen so far:
		% \begin{equation}
			%   \tilde{I}_t(x,y) = \min_{k \leq t}\, I_k(x,y).
			%   \label{eq:min_history}
			% \end{equation}
		% This enforces monotonic darkening in time: since damage only accumulates and should only reduce transmitted light, any transient brightening (illumination flashes, sensor noise) is suppressed \cite{benderEffectVariableAmplitude2021}.
		
		% \item \textbf{Reference normalization.} The history-clamped frame $\tilde{I}_t$ is divided pixel-wise by a reference image $R_t$ to cancel features already present before frame $t$ and highlight newly appeared damage:
		% \begin{equation}
			%   N_t(x,y) = \operatorname{clip}\!\left(\frac{\tilde{I}_t(x,y)}{R_t(x,y)},\; 0,\; 1\right),
			%   \label{eq:normalization}
			% \end{equation}
		% where the denominator is bounded below by a small constant to avoid division by zero. Two reference strategies are available. In \emph{static} mode, $R_t$ is fixed to a single undamaged frame captured before damage onset, so normalization always measures deviation from the pristine state. In \emph{rolling-median} mode, $R_t$ is computed as the pixel-wise median of a sliding window of $w$ frames ending $s$ frames before $t$:
		% \begin{equation}
			%   R_t(x,y) = \operatorname{median}\!\left\{\tilde{I}_{t-s-w}(x,y),\,\ldots,\, \tilde{I}_{t-s-1}(x,y)\right\},
			%   \label{eq:rolling_median}
			% \end{equation}
		% with parameters $w$ (\texttt{reference\_window}) and $s$ (\texttt{reference\_skip}). The rolling-median reference adapts gradually to slow background changes (e.g.\ lens contamination or illumination drift during long fatigue tests), whereas the static reference yields a stronger contrast signal. The result $N_t$ is rescaled to the $[0,255]$ integer range before further processing \cite{benderEffectVariableAmplitude2021, olesenMatrixCraCSAutomatedTracking2024}.
		
		% \item \textbf{Max/min filtering.} A morphological open by reconstruction (sequential dilation then erosion) with a rectangular structuring element of size $w_y \times w_x$ regularizes local thin structures such as crack lines that could otherwise be misinterpreted as delamination:
		% \begin{align}
			%   F_t^+(x,y) &= \max_{(u,v) \in \mathcal{W}}\, N_t(x+u,\, y+v), \label{eq:max_filter}\\
			%   F_t(x,y)   &= \min_{(u,v) \in \mathcal{W}}\, F_t^+(x+u,\, y+v), \label{eq:min_filter}
			% \end{align}
		% where $\mathcal{W} = \{-\lfloor w_y/2\rfloor,\ldots,\lfloor w_y/2\rfloor\} \times \{-\lfloor w_x/2\rfloor,\ldots,\lfloor w_x/2\rfloor\}$.
		
		% \item \textbf{Sharpening and Gaussian smoothing.} An unsharp-mask filter enhances local contrast around delamination boundaries:
		% \begin{equation}
			%   S_t = (1+\alpha)\,F_t - \alpha\,(G_{\sigma_1} * F_t),
			%   \label{eq:unsharp}
			% \end{equation}
		% with gain $\alpha = 2$ and $\sigma_1 = \bar{w}_c$ (mean crack width in pixels, estimated from the crack catalogue). The high-pass boost widens the apparent intensity contrast between delaminated and intact regions. A subsequent anisotropic Gaussian filter suppresses residual noise while preserving the elongated delamination front:
		% \begin{equation}
			%   \hat{S}_t = G_{(\sigma_{2,y},\,\sigma_{2,x})} * S_t,
			%   \label{eq:gauss_smooth}
			% \end{equation}
		% with default $(\sigma_{2,y}, \sigma_{2,x}) = (0.5, 15.0)$ pixels.
		
		% \item \textbf{Intensity scaling.} To decouple the threshold from frame-wise contrast variations, intensities are linearly mapped to $[0,1]$ using percentile-based bounds:
		% \begin{equation}
			%   \hat{I}_t(x,y) = \operatorname{clip}\!\left(\frac{\hat{S}_t(x,y) - p_{\alpha_l}}{p_{\alpha_u} - p_{\alpha_l}},\; 0,\; 1\right),
			%   \label{eq:scaling}
			% \end{equation}
		% where $p_{\alpha_l}$ and $p_{\alpha_u}$ are the $\alpha_l$th and $\alpha_u$th percentiles of $\hat{S}_t$ (defaults: $\alpha_l = 10$, $\alpha_u = 99$). This ensures that the threshold behaviour is stable across frames regardless of absolute brightness changes \cite{bakMethodAutomatedDigital2020}.
		
		% \item \textbf{Thresholding.} A one-dimensional two-cluster k-means is fitted to the pixel values of $\hat{I}_t$, initialized from the minimum and maximum values and iterated until convergence:
		% \begin{equation}
			%   \tau^* = \tfrac{1}{2}(\mu_0 + \mu_1),
			%   \label{eq:kmeans_threshold}
			% \end{equation}
		% where $\mu_0 < \mu_1$ are the converged cluster centroids. If the algorithm does not converge or if both clusters collapse to the same value, the Otsu threshold \cite{otsuThresholdSelectionMethod1979} is used as a fallback. The binary candidate mask is then:
		% \begin{equation}
			%   B_t(x,y) = \mathbb{1}\!\left[\hat{I}_t(x,y) < \tau^*\right] \;\cap\; \mathbb{1}\!\left[\hat{S}_t(x,y) \leq \eta \cdot 255\right],
			%   \label{eq:binary_mask}
			% \end{equation}
		% where $\eta \in (0,1)$ is a hard-floor parameter (\texttt{hard\_floor}, default $0.90$) that additionally suppresses pixels whose smoothed intensity is too bright to be considered damaged, regardless of the threshold \cite{bakMethodAutomatedDigital2020, olesenMatrixCraCSAutomatedTracking2024}.
		
		% \item \textbf{Morphological closing.} Small holes and narrow gaps in $B_t$ caused by image texture or thin undamaged regions within a delaminated zone are filled via morphological closing with a disk-shaped structuring element of radius $r$:
		% \begin{equation}
			%   \bar{B}_t = B_t \;\bullet\; D_r,
			%   \label{eq:closing}
			% \end{equation}
		% where $\bullet$ denotes the closing operator and $D_r = \{(u,v): u^2+v^2 \leq r^2\}$ \cite{bakMethodAutomatedDigital2020}.
		
		% \item \textbf{Temporal accumulation.} The accumulated mask is updated as the running union of all per-frame candidates, enforcing physically consistent monotonic growth:
		% \begin{equation}
			%   M_t = M_{t-1} \;\cup\; \bar{B}_t, \qquad M_0 = \emptyset.
			%   \label{eq:temporal_accumulation}
			% \end{equation}
		% This suppresses frame-to-frame flicker and ensures that once a region is detected as delaminated it is not later reverted.
		
		
		% \item \textbf{Max/min filtering.} A grayscale morphological closing (dilation followed by erosion) with a rectangular structuring element of size $w_y \times w_x$ regularizes local thin structures such as crack lines that could otherwise be misinterpreted as delamination:
		% \begin{align}
			%   F_t^+(x,y) &= \max_{(u,v) \in \mathcal{W}}\, N_t(x+u,\, y+v), \label{eq:max_filter}\\
			%   F_t(x,y)   &= \min_{(u,v) \in \mathcal{W}}\, F_t^+(x+u,\, y+v), \label{eq:min_filter}
			% \end{align}
		% where $\mathcal{W} = \{-\lfloor w_y/2\rfloor,\ldots,\lfloor w_y/2\rfloor\} \times \{-\lfloor w_x/2\rfloor,\ldots,\lfloor w_x/2\rfloor\}$. The first stage, $F_t^+$, propagates local bright values into narrow dark valleys, while the second stage restores the background envelope without re-introducing dark structures that are smaller than the structuring support. In practice, this suppresses dark crack lines more reliably than using either stage alone.
		
		% For parameterization, $w_x$ should be tied to the measured mean crack width $\bar{w}_c$ through a multiplier, e.g. $w_x = k\,\bar{w}_c$ with $k>1$ (typically $k\in[1.5,3]$), while keeping $w_y$ small. Using $w_x\approx\bar{w}_c$ is often insufficient in experimental data because apparent crack width varies with blur, local orientation, illumination, and pixel discretization.
		
		
		% \end{enumerate}
	
	
	
	
	\section*{Declaration of competing interest}
	The authors declare no competing financial interests or personal relationships.
	
	
	\section*{Acknowledgements}
	The authors would like to thank the Institute of Materials Science and Testing of Polymers, Department of Polymer Engineering and Science, TU Leoben, for supporting this work by providing the examples pictures presented in this paper.
	
	
	\bibliographystyle{elsarticle-num}
	\bibliography{bibliography}
	
	
	
	
\end{document}

